% Generated draft: requires mathematical and source review.
\documentclass[11pt]{article}
\usepackage[margin=25mm]{geometry}
\usepackage{fontspec}
\setmainfont{Latin Modern Roman}
\usepackage{amsmath,amssymb,mathtools,unicode-math}
\setmathfont{Latin Modern Math}
\usepackage{enumitem,needspace,xurl,hyperref}
\setlength{\parindent}{0pt}
\setlength{\parskip}{5pt}
\setlength{\emergencystretch}{3em}
\newcommand{\N}{\mathbb{N}}
\newcommand{\Z}{\mathbb{Z}}
\newcommand{\Q}{\mathbb{Q}}
\newcommand{\R}{\mathbb{R}}
\newcommand{\C}{\mathbb{C}}
\newcommand{\F}{\mathbb{F}}
\newcommand{\A}{\mathbb{A}}
\newcommand{\PP}{\mathbb{P}}
\newcommand{\E}{\mathbb{E}}
\newcommand{\eps}{\varepsilon}
\newcommand{\dd}{\,\mathrm{d}}
\begin{document}
\section*{UnsolvedMath ID: 30001778}
\textbf{Source problem code:} catalogue entry
\subsection{Definitions and mathematical statement}
For a label set \(I\), a partition of \(I\) is a collection of pairwise disjoint nonempty subsets, called blocks, whose union is \(I\). A partition \(\pi'\) is coarser than \(\pi\) if every block of \(\pi\) is contained in a block of \(\pi'\). In a partition-valued formulation, a coalescent process \((\Pi_t)_{t\in T}\) evolves by merging blocks: whenever \(s\leq t\), \(\Pi_t\) is coarser than \(\Pi_s\), almost surely. Here \(T\) denotes the time index set; the source does not specify it or the label set \(I\). Additional process assumptions, including Markov and path regularity requirements, are not specified by the supplied statement.

For a permutation \(\sigma\) of \(I\), write \(\sigma\pi=\{\sigma(B):B\in\pi\}\). Exchangeability means invariance of the process law under relabeling; in particular, for every positive integer \(k\), every \(t_1,\ldots,t_k\in T\), and every permutation \(\sigma\) moving only finitely many labels,
\[
 (\sigma\Pi_{t_1},\ldots,\sigma\Pi_{t_k})
 \overset{\mathrm d}{=}
 (\Pi_{t_1},\ldots,\Pi_{t_k}).
\]
These definitions explain the terminology without selecting a finite or infinite label domain for the requested example.

\textbf{Source problem.} Find another exchangeable coalescent that can be represented by cutting down a combinatorial tree.

The requested construction would have to specify a combinatorial tree model, a cutting procedure, and a rule obtaining the coalescent partitions from that procedure. The supplied record does not identify the previously known examples relative to which the example must be ``another.'' It also does not specify the tree distribution, whether trees are finite or infinite, which objects are cut, how labels are assigned to blocks, or whether the representation involves reversing time. Consequently, the displayed source problem is a construction request whose intended admissible class cannot be stated as a fully determined existential assertion from the supplied evidence. No particular representation convention or excluded example is imposed here.
\subsection{Short English statement}
Find a further exchangeable process of merging partition blocks that admits a representation through cutting a combinatorial tree. The supplied statement leaves the known examples to be excluded, the tree model, and the precise meaning of the representation unspecified. This is a statement-only draft requiring human mathematical review.
\subsection{Sources}
\begin{itemize}
\item Local UnsolvedMath registry: \texttt{data/unsolvedmath/problems.json}, permanent ID \(30001778\), source problem code \texttt{catalogue entry}, title \emph{Tree-Cutting Representations of Exchangeable Coalescents}. The supplied original, clean, and statement fields contain the same construction request.
\item Supplied report citation: \emph{Mini-Workshop: Random Trees, Information and Algorithms} (2011), \emph{Oberwolfach Reports}, pp.\ 1241--1286, DOI \texttt{10.4171/owr/2011/23}. Source URL and supplied literature-source link, labeled ``original OWR report'': \url{https://doi.org/10.4171/owr/2011/23}. No individual proposer is identified in the record; the supplied proposed year is 2011.
\item The registry marks the statement as ``exact'' and supplies a verification note asserting comparison with the official report and normalization of notation, context, and question form. This is supplied provenance, not independent verification in this draft.
\item Literature assessment dated \textbf{2026-08-21}: the supplied status is ``open.'' The supplied assessment says that the extracted text is only an ``Open problem'' label with no mathematical statement and was classified as open under the queue policy. This assessment conflicts with the construction request present in the supplied statement fields. It is supplied dated curation, not independently verified literature assessment or confirmation of current mathematical status. The background likewise calls for independent verification of time-sensitive status and cited claims.
\end{itemize}
\end{document}
